\documentclass[conference,a4paper]{IEEEtran}
\usepackage[utf8]{inputenc}
\usepackage{booktabs}
\usepackage{url}

\begin{document}

\title{Learning Individual Comfort: Personalized Proxemic Greeting\\
for a Simulated Social Robot via Reinforcement Learning}

\author{\IEEEauthorblockN{Eman Saleh, Nour Gaser}
\IEEEauthorblockA{Department of Computer Science, German International University (GIU), Cairo, Egypt\\
eman.khalil@alumni2024.giu-uni.de, nouraldin.algendi@student.giu-uni.de}}

\maketitle

\section{Problem Addressed}

Social robots optimised for task completion routinely ignore that individuals
differ: the interpersonal distance and greeting intensity one person finds
welcoming, another finds intrusive. Reinforcement learning has been shown to adapt
robot behaviour to users \cite{wang2025,tozadore2018,gonzalez2024,zedda2023}, but
those evaluations are small ($N{=}8$--$28$), hardware-bound, and rarely separate
\emph{adaptation} --- learning a good behaviour --- from \emph{personalization}
--- learning a \emph{different} behaviour for each person.

This work addresses that distinction: \emph{does a social robot actually need to
personalize, or is one well-trained policy enough for everyone?} Personalization is
expensive --- it requires per-user interaction data --- so answering this demands
comparison against a \emph{trained generalist}, not merely a scripted robot.

\section{Approach}

We implement greeting/proxemics as a discrete MDP in simulation. State is the
sensed distance bin (Hall's proxemic zones) crossed with commanded greeting
intensity ($4\times3=12$ states); actions are $\{$approach, withdraw, intensify,
soften, hold$\}$. Reward uses social signals only --- $+1$ approval, $-2$ if the
person disengages, $-0.05$ per step --- so the robot never observes the preference
it must satisfy. Learning is tabular Q-learning: interpretable, and 12 states need
no function approximation and no GPU.

Individual difference is modelled by a \texttt{UserProfile}: three simulated people
(\emph{reserved}, \emph{neutral}, \emph{sociable}) with \emph{different hidden}
comfortable distances and preferred intensities, placed in mutually exclusive
distance bins so a policy optimal for one is wrong for another. We compare five
conditions: scripted; one policy trained across the mixed population; the
strongest \emph{coherent} single policy; personalized per user; and warm-started
from other users then adapted (leave-one-out).

All work is in simulation: without access to a physical humanoid, hardware trials
and a subjective human study are explicitly \textbf{out of scope} and remain future
work. The learning code addresses an abstract robot interface rather than any
specific platform, so the identical loop would run against hardware unchanged.

\section{Expected Outcome}

The implementation is complete; preliminary results follow (8 seeds, 6000
episodes, mean reward per interaction):

\vspace{-0.8em}
\begin{table}[h]
\centering
\scriptsize
\begin{tabular}{@{}lccc@{}}
\toprule
\textbf{Condition} & \textbf{Reserved} & \textbf{Neutral} & \textbf{Sociable} \\
\midrule
Scripted            & $-1.26$          & $2.03$           & $-0.75$ \\
One policy, mixed   & $6.45$           & $-0.75$          & $-0.75$ \\
Best single policy  & $\mathbf{11.81}$ & $-0.75$          & $-0.75$ \\
Personalized        & $\mathbf{11.83}$ & $\mathbf{10.31}$ & $\mathbf{8.34}$ \\
\bottomrule
\end{tabular}
\end{table}
\vspace{-1.2em}

The result is unambiguous: the best single non-personalized policy is optimal for
\emph{one} user ($11.81$) and worse than doing nothing for the other two ($-0.75$),
while personalization gains $+11.06$ and $+9.09$ there. It is not a refinement; it
is the difference between a working greeting and a failed one.

We also report a \textbf{negative result}: warm-starting from other users did not
help, merely matching from-scratch learning at equal budget and faring worse at a
small one. We attribute this to a boundary condition on the transfer idea in PbARL
\cite{wang2025}, which works because task competence and user preference are
\emph{separable} --- the robot must still complete the task, only differently. In
proxemics the task \emph{is} the preference, so another person's optimum carries no
information about this one's. This is a useful limit rather than a failure, and the
result most worth discussing.

The demonstration presents learning curves, the per-user learned policies side by
side, and the comparison above, executed live. The implementation is publicly
available\footnote{\url{https://github.com/nourgaser/nao-adaptive-rl}} with a test
suite; every reported number is reproducible from a fixed seed.

\vspace{-0.5em}
\begin{thebibliography}{9}
\scriptsize
\setlength{\itemsep}{0pt}
\bibitem{wang2025} R. Wang \textit{et al.}, ``Personalization in HRI through preference-based action representation learning,'' \textit{IEEE ICRA}, 2025.
\bibitem{tozadore2018} D. C. Tozadore \textit{et al.}, ``Towards adaptation and personalization in task based on HRI,'' \textit{LARS/SBR/WRE}, 2018.
\bibitem{gonzalez2024} A. Gonzalez-Santocildes \textit{et al.}, ``Adaptive robot behavior based on human comfort using RL,'' \textit{IEEE Access}, 12, 2024.
\bibitem{zedda2023} E. Zedda \textit{et al.}, ``Adaptive behaviour-based strategy for SARs with MCI older adults,'' \textit{arXiv:2305.01492}, 2023.
\end{thebibliography}

\end{document}
